\documentclass[11pt,reqno]{amsart}
\usepackage[margin=1.15in]{geometry}
\usepackage{amsmath,amssymb,amsthm,mathtools}
\usepackage{booktabs}
\usepackage[colorlinks=true,linkcolor=blue,citecolor=blue,urlcolor=blue]{hyperref}

\theoremstyle{plain}
\newtheorem{theorem}{Theorem}[section]
\newtheorem{proposition}[theorem]{Proposition}
\newtheorem{lemma}[theorem]{Lemma}
\newtheorem{corollary}[theorem]{Corollary}
\theoremstyle{definition}
\newtheorem{computation}[theorem]{Computation}
\newtheorem{remark}[theorem]{Remark}
\newtheorem{example}[theorem]{Example}

\DeclareMathOperator{\Pic}{Pic}
\DeclareMathOperator{\Jac}{Jac}
\DeclareMathOperator{\divisor}{div}
\DeclareMathOperator{\ord}{ord}
\newcommand{\Q}{\mathbb{Q}}
\newcommand{\Z}{\mathbb{Z}}
\newcommand{\C}{\mathbb{C}}
\newcommand{\HH}{\mathcal{H}}

\title{The character of a Borcherds form on a quotient of positive genus,\\ and the reduction type of a genus-one quotient}
\author{}
\date{}

\begin{document}
\maketitle

\section{Purpose}

Guo and Yang \cite{GY} construct equations of Shimura curves $X_0^D(N)/W$ from Borcherds
forms whose values at CM points are given by Schofer's formula. When the quotient has genus
zero the construction is unconditional. For a quotient of positive genus it rests on their
Assumption~40, which says that a Borcherds form satisfying the criterion of their Lemma~7 has
trivial character. This note records what that assumption amounts to, when it holds, a case in
which it fails, how to do without it, and a way to determine which curve of an isogeny class a
genus-one quotient is that uses nothing of the Borcherds construction.

Throughout, $X = X_0^D(N)/W$ is a quotient of a Shimura curve by a group $W$ of Atkin--Lehner
involutions, $\Gamma$ is the Fuchsian group with $X(\C) = \Gamma\backslash\HH$, and $\psi$ is a
Borcherds form on $X$ of weight $0$ attached to a weakly holomorphic form $F$ with integral
principal part, so that $\psi$ is a meromorphic function on $\HH$ with
$\psi(\gamma z) = \chi(\gamma)\,\psi(z)$ for a character $\chi$ of $\Gamma$, and
$D = \divisor(\psi)$ is a $\Q$-rational divisor of degree $0$ on $X$ supported on CM points. At an
elliptic point of order $k$ the order of $\psi$ on $X$ is its order on $\HH$ divided by $k$
(\cite{GY}, Lemma~7), so it is the integrality of the principal part of $F$, through the parity
condition recalled in \S2, that makes $D$ a divisor with integral multiplicities and the word
``principal'' meaningful for it. The character $\chi$ has finite order, say $n$: this is part of
Borcherds' theorem on the Heegner divisors that are principal (\cite{Bo}, proof of Theorem~4.1),
whose published version omitted the finite-order condition from the definition of a principal
Heegner divisor; the corrected version \cite{BoCorr} restores it and fills the gap.

\section{What Assumption 40 says}

The group $\Gamma$ is generated by hyperbolic elements $\alpha_1, \beta_1, \dots, \alpha_g,
\beta_g$ and elliptic elements $\gamma_1, \dots, \gamma_r$ of orders $k_1, \dots, k_r$, subject to
$\prod_j [\alpha_j, \beta_j] \prod_i \gamma_i = 1$ and $\gamma_i^{k_i} = 1$. Lemma~7 of \cite{GY}
(Lemma~8 of \cite{Yang}) says: if $c_0(0)$ is even, then $\chi(\gamma_i) = 1$ for \emph{every} $i$ if
and only if, at each elliptic point of order $k_i \neq 3$, the order of $\psi$ has the parity of
$c_0(0)/2$. The points of order $3$ need no condition of their own: a character kills commutators,
so the relation $\prod_j [\alpha_j, \beta_j] \prod_i \gamma_i = 1$ gives $\prod_i \chi(\gamma_i) = 1$,
and $\chi(\gamma_i)$ has order dividing $k_i$, which for $k_i = 3$ is coprime to the orders at the
other elliptic points. This is a condition on the divisor $D$ alone. When $g = 0$ there are no
hyperbolic generators, so it decides $\chi$; this is why the equations in the main table of
\cite{GY} are unconditional. Assumption~40 is the statement that the same parity condition
decides $\chi$ when $g > 0$; as Remark~41 of \cite{GY} says, what is assumed is
$\chi(\alpha_j) = \chi(\beta_j) = 1$.

\section{The character is trivial exactly when the divisor is principal}

\begin{proposition}\label{prop:principal}
Suppose the parity condition of Lemma~7 holds, so that $\chi$ is trivial on the elliptic
generators. Then the following are equivalent:
\begin{enumerate}
\item $\chi$ is trivial, i.e.\ $\psi$ is a function on $X$;
\item $D$ is a principal divisor on $X$;
\item the class of $D$ vanishes in $\Jac(X)(\Q)_{\mathrm{tors}}$.
\end{enumerate}
In general the order of $\chi$ is the order of the class of $D$ in $\Jac(X)(\Q)$.
\end{proposition}

\begin{proof}
Since $\chi$ has order $n$, $\psi^n$ is invariant under $\Gamma$, hence a function on $X$ with
divisor $nD$; so the class of $D$ is torsion, and it is a $\Q$-rational class since $D$ is
$\Q$-rational. If $D = \divisor(h)$ for a function $h$ on $X$, then $u = \psi/h$ is a
holomorphic function on $\HH$ without zeros or poles with $u(\gamma z) = \chi(\gamma) u(z)$.
Since $|\chi| = 1$, $\log|u|$ is a $\Gamma$-invariant harmonic function on $\HH$, hence a
harmonic function on the compact Riemann surface $X$ (bounded near the elliptic points, so it
extends across them), hence constant. A holomorphic function of constant modulus is constant, so
$u$ is constant and $\chi = 1$. Conversely if $\chi = 1$ then $D = \divisor(\psi)$ is principal. The last sentence follows by applying the
argument to $\psi^m$ for $m \mid n$.
\end{proof}

\begin{corollary}\label{cor:torsionfree}
If the degree-zero CM divisors of $X$ generate a subgroup of $\Jac(X)(\Q)$ with no torsion, in
particular if $X$ has genus one and $E = \Jac(X)$ has $E(\Q)_{\mathrm{tors}} = 0$, then every
Borcherds form on $X$ satisfying the parity condition has trivial character: Assumption~40 holds
for every $F$. In general the order of $\chi$ divides the exponent of $\Jac(X)(\Q)_{\mathrm{tors}}$.
\end{corollary}

So Assumption~40 is unconditional for $X_0^{142}(1)^*$, $X_0^{158}(1)^*$ and $X_0^{38}(5)^*$, whose
Jacobians $142a1$, $158a1$ and $190a1$ have no rational torsion, and it cannot be decided from $F$
alone when there is rational torsion: the class of $D$ in $\Jac(X)(\Q)_{\mathrm{tors}}$ is not a
function of the discriminants and multiplicities in $D$.

\begin{remark}[How a failure shows]\label{rem:failure}
Schofer's formula gives $\log|\psi|$ at a CM point, and $|\psi|$ is well defined on $X$ whatever
$\chi$ is, since $|\chi| = 1$. The statement to use is about \emph{ratios}, which avoids any
assumption on the normalising constant. Let $Q_1, Q_2$ be CM points that are rational on $X$ and
not in the support of $D$. If $\chi = 1$ then $\psi = c\,h$ with $h \in \Q(X)$ and $c$ a constant,
so $|\psi(Q_1)/\psi(Q_2)| = |h(Q_1)/h(Q_2)|$ is the absolute value of a rational number, whatever
$c$ is. If $\chi$ has order $n$ then the same applies to $\psi^n$, so that ratio is an $n$-th root
of a rational number, and for $n = 2$ it is in general the square root of a rational number that is
not a square. A ratio such as $2\cdot 7^{3/2}\cdot 17^{1/2}$ is therefore the signature of a
character of order $2$, and it says nothing about the constant.
\end{remark}

\begin{example}[Assumption 40 appears to fail at $X_0^{14}(17)^*$]\label{ex:14-17}
The star curve $X = X_0^{14}(17)/W_{14,17}$ has genus one and lies in the isogeny class $238a$,
whose two curves have $E(\Q) \cong \Z/2 \times \Z$, so Corollary~\ref{cor:torsionfree} does not
apply and the assumption is a genuine hypothesis here. A run of the driver that builds models over a
genus-one star reported a Borcherds form on $X$, satisfying the parity condition of Lemma~7, whose
value at the CM point of discriminant $-4$ is $2\cdot 7^{3/2}\cdot 17^{1/2}$ (the exponents $3/2$ and
$1/2$ sit at the primes $7 \mid D$ and $17 \mid N$). That point is rational on $X$ and is the only CM
point of discriminant $-4$ there, so Remark~\ref{rem:failure} applies to it and the character of that
form has order $2$: the class of its divisor is the nonzero element of $E(\Q)[2]$. We record this as
an observation from one run rather than as a theorem: the value, the divisor of the form and the
orders at the elliptic points are in that driver's log, which lives with the driver and not with
this note, so the example is not reproducible from the files beside it.
\end{example}

\section{Doing without the assumption}

Let $e$ be the exponent of $\Jac(X)(\Q)_{\mathrm{tors}}$ (for a genus-one quotient, of
$E(\Q)_{\mathrm{tors}}$, taken over the curves of the isogeny class when the curve is not yet
known). By Proposition~\ref{prop:principal}, $\psi^e$ is a function on $X$ for every Borcherds
form $\psi$ satisfying the parity condition, with divisor $eD$, and Schofer's formula gives
$|\psi^e| = |\psi|^e$ exactly. Building the coordinates of \cite{GY}, \S4.3 from $\psi^e$ instead of
$\psi$ therefore needs no assumption: the coordinates have degree $e$ times what they would
otherwise have, the relation between them has correspondingly higher bidegree, and nothing else
changes. When $E(\Q)$ is torsion-free, $e = 1$. The repository's driver takes this route on
quotients with rational torsion.

\section{The reduction type of a genus-one quotient from the dual graph}

The curve inside the isogeny class can be read off the Čerednik--Drinfeld model, with nothing of the
Borcherds construction involved. This is how González and Rotger identified the Jacobians of the
genus-one Shimura curves and of their non-elliptic genus-one Atkin--Lehner quotients (\cite{GR},
the table in the proof of Theorem~3.3 and Table~2 of Theorem~4.2); what follows is their method,
with the metric dual graph of the repository supplying the thicknesses.

Let $p \mid D$. Over $\Z_p^{\mathrm{nr}}$ the Čerednik--Drinfeld model of $X$ has special fibre a
union of rational curves meeting at nodes, and its dual graph with each edge weighted by the
thickness of the node is the quotient by $W$ of the dual graph of $X_0^D(N)$, which is the quotient
of the Bruhat--Tits tree at $p$; the thickness of a node is the order of the stabiliser of the
corresponding edge in the group acting (\cite{Ku}, Prop.~3-2; see also \cite{BC}). The lengths
computed by \texttt{DualGraphMetricQuotientFromData} are stabiliser orders in the group containing
the centre $\{\pm 1\}$, hence twice the thickness.

\begin{proposition}\label{prop:kodaira}
Let $X$ have genus one with Jacobian $E$, and let $p \mid D$. Then $E$ has multiplicative reduction
at $p$ of Kodaira type $I_n$, where $n$ is the sum of the thicknesses of the nodes lying on the
unique cycle of the dual graph, equivalently half the total length of the metric dual graph after
its trees are contracted.
\end{proposition}

\begin{proof}
First, $X$ and $E$ have the same reduction type over $\Z_p^{\mathrm{nr}}$. The special fibre of the
Čerednik--Drinfeld model is a curve over the algebraically closed residue field, so it has a smooth
rational point, which lifts by smoothness; thus $X(\Q_p^{\mathrm{nr}}) \neq \emptyset$ and $X$ is
isomorphic to $E$ over $\Q_p^{\mathrm{nr}}$. Their minimal regular models over
$\Z_p^{\mathrm{nr}}$ therefore agree, and the index of $I_n$ does not change under unramified base
change, so it may be computed from either.

Now pass from the Čerednik--Drinfeld model to the minimal regular one. Resolving a node of
thickness $t$ replaces it by a chain of $t - 1$ rational curves (\cite{Liu}, \S10.3), and the
components meeting the rest of the fibre in a single point are then contracted, which is exactly
the contraction of the trees of the dual graph: what remains is a cycle of rational curves, the
special fibre of a curve of arithmetic genus one with multiplicative reduction. For type $I_n$ that
cycle has $n$ components (\cite{Si}, IV.9), and the number of components after the resolution is the
sum over the nodes on the cycle of their thicknesses. Halving the metric lengths accounts for the
centre.
\end{proof}

\begin{remark}[Scope]\label{rem:scope}
The method sees only the primes dividing $D$, so two curves of an isogeny class differing only at a
prime of $N$ would not be separated by it; that does not occur for the classes below. What it
identifies is the Jacobian, which is the curve itself whenever the quotient has a rational point, as
in every case here.
\end{remark}

\begin{computation}
\texttt{KodairaIndexFromDualGraph} computes $n$; \texttt{tests/KodairaFromDualGraph.m} checks it.
Calibration on the three stars whose curve is known (the newform's isogeny class has a single
curve): all six indices agree with Cremona's tables. Identification of the two stars known only up
to isogeny from the Borcherds construction:
\begin{center}
\begin{tabular}{llll}
\toprule
quotient & dual graph & curves of the class & curve \\
\midrule
$X_0^{142}(1)^*$ & $I_9$ at $2$, $I_1$ at $71$ & $142a1$: $I_9, I_1$ & $142a1$ \\
$X_0^{158}(1)^*$ & $I_8$ at $2$, $I_1$ at $79$ & $158a1$: $I_8, I_1$ & $158a1$ \\
$X_0^{38}(5)^*$ & $I_{11}$ at $2$, $I_1$ at $19$ & $190a1$: $I_{11}, I_1$ & $190a1$ \\
$X_0^{14}(17)^*$ & $I_7$ at $2$, $I_4$ at $7$ & $238a1$: $I_{14}, I_2$; $238a2$: $I_7, I_4$ & $238a2$ \\
$X_0^{94}(3)^*$ & $I_8$ at $2$, $I_1$ at $47$ & $282b1$: $I_8, I_1$; $282b2$: $I_4, I_2$ & $282b1$ \\
\bottomrule
\end{tabular}
\end{center}
Both identifications agree with the curves the Borcherds construction produced under
Assumption~40, including at $X_0^{14}(17)^*$ where the assumption appears to fail for some forms and
the model was built from others. The outside check is on the $29$ genus-one curves of \cite{GR}:
the six curves $X_0^D(1)$ for $D = 14, 15, 21, 33, 34, 46$, their quotients by one involution, and
the seventeen quotients of Table~2 from $(39,13)$ to $(462,154)$. In every case the indices at the
primes of $D$ single out one curve of the isogeny class, and it is the curve named there. That run
is \texttt{tests/KodairaGonzalezRotger.m}, which takes the isogeny class from the trace formula
rather than from a label -- the Jacobian is $D$-new, so its conductor is $D$ and its traces of
Frobenius identify the class -- and compares the published label only where this repository has
transcribed it. It includes edge stabilisers of order $2$ at $p = 2$ and of order $3$ at $p = 3$, and
quotients by $w_p$ itself. For the two quotients of this note the heights of the branch points on
the quotient curve, computed by the Heegner-height route of the repository, gave the same answer;
that is a second agreement, not an independent proof, since we have no argument that the predicted
height always fails to be a lattice height on the wrong curve of a class.
\end{computation}

\begin{thebibliography}{9}
\bibitem{GY} J.-W.~Guo and Y.~Yang, Equations of hyperelliptic Shimura curves, Compositio Math.\
153 (2017), 1–40.
\bibitem{Yang} Y.~Yang, Special values of hypergeometric functions and periods of CM elliptic
curves, arXiv:1503.07971 (2015); Lemma~8.
\bibitem{Bo} R.~E.~Borcherds, The Gross--Kohnen--Zagier theorem in higher dimensions, Duke Math.\ J.\
97 (1999), 219--233.
\bibitem{BoCorr} R.~E.~Borcherds, Correction to ``The Gross--Kohnen--Zagier theorem in higher
dimensions'', Duke Math.\ J.\ 105 (2000), 183--184; corrected version arXiv:alg-geom/9710002v3.
\bibitem{GR} J.~González and V.~Rotger, Non-elliptic Shimura curves of genus one, J.\ Math.\ Soc.\
Japan 58 (2006), 927--948; arXiv:math/0612732.
\bibitem{Ku} A.~Kurihara, On some examples of equations defining Shimura curves and the Mumford
uniformization, J.\ Fac.\ Sci.\ Univ.\ Tokyo 25 (1979), 277--300; Prop.~3-2.
\bibitem{BC} J.-F.~Boutot and H.~Carayol, Uniformisation $p$-adique des courbes de Shimura,
Astérisque 196--197 (1991), 45--158.
\bibitem{Liu} Q.~Liu, Algebraic Geometry and Arithmetic Curves, Oxford, 2002; \S10.3.
\bibitem{Si} J.~H.~Silverman, Advanced Topics in the Arithmetic of Elliptic Curves, Springer, 1994;
IV.9.
\end{thebibliography}

\end{document}
